\section{Hierarchy theorem}

\subsection{Codifica di una macchina di Turing}
$\underline{M}$ é una stringa binaria che codifica una macchina di Turing $M$. \\
I simboli dell'alfabeto, gli stati e gli spostamenti della testina sono codificati con un numero naturale.
La relazione (o funzione) di transizione é codificata da una quintupla di numeri naturali.

\[\underline{M} = (\underline{\Sigma}, \underline{Q}, \underline{\Delta})\]

\begin{itemize}
    \item $\sigma \in \Sigma \quad \rightsquigarrow \quad \underline{\sigma} \in \{1, \dots, |\Sigma|\}$
    \item $q \in Q \quad \rightsquigarrow \quad \underline{q} \in \{|\Sigma| + 1, \dots, |\Sigma| + |Q|\}$
    \begin{itemize}
        \item $s     \quad \rightsquigarrow \quad \underline{s}     = |\Sigma| + 1$
        \item $"yes" \quad \rightsquigarrow \quad "\underline{yes}" = |\Sigma| + |Q| - 2$
        \item $"no"  \quad \rightsquigarrow \quad "\underline{no}"  = |\Sigma| + |Q| - 1$
        \item $h     \quad \rightsquigarrow \quad \underline{h}     = |\Sigma| + |Q|$
    \end{itemize}
    \item $(q, [\sigma], \quad q', [\sigma'], [dir]) \in \Delta \rightsquigarrow (\underline{q}, [\underline{\sigma}], \quad \underline{q'}, [\underline{\sigma'}], [\underline{dir}]) \in \underline{\Delta}$
    \subitem $dir \in \{ \leftarrow, \rightarrow, -\}$, dove:
    \begin{itemize}
        \item $"\leftarrow" \quad \rightsquigarrow \quad |\Sigma| + |Q| + 1$
        \item $"\rightarrow" \quad \rightsquigarrow \quad |\Sigma| + |Q| + 2$
        \item $"- " \quad \rightsquigarrow \quad |\Sigma| + |Q| + 3$
    \end{itemize}
\end{itemize}

\subsubsection{Lunghezza della codifica}
Servono $|\Sigma| + |Q| + 3$ valori per codificare tutti i simboli, stati o spostamenti. 
La stringa binaria che ne codifica uno é lunga \footnote{Sulle slide é riportato il caso particolare con $k = 1$}:
\[\ell = \lceil \log_2(|\Sigma| + |Q| + 3) \rceil\]
Di conseguenza, la codifica di una transizione in $\underline{\Delta}$ é lunga:
\begin{itemize}
    \item $2$ per le parentesi
    \item $\ell$ per lo stato di partenza
    \item $k \ell$ per i simboli sotto le testine
    \item $\ell$ per lo stato di arrivo
    \item $k \ell$ per i simboli da scrivere
    \item $k \ell$ per gli spostamenti delle testine
\end{itemize}
\[2 + 2\ell + 3k \ell\]
E quindi $\underline{\Delta}$ é lungo:
\[|\underline{\Delta}| = |\Delta| \cdot (2 + 2\ell + 3k \ell)\]

Ma se $|\underline{\Delta}|$ é logaritmico in $N$, possiamo concludere, pessimisticamente, che: 
\[|\Delta| = O((\log(N))^2)\]

\subsection{Macchina di Turing Universale}
\begin{itemize}
    \item $U$: Macchina di Turing Universale
    \item $X = (\underline{M}, x)$: Input di $U$
    \item $N = |X| = |\Sigma| + |Q| + |\Delta| + n $
\end{itemize}
Se $M$ ha $k$ nastri, possiamo imporre che $U$ abbia $k+1$ nastri (ne serve uno per mantenere $X$ e altri $k$ per mantenere la configurazione di $M$), 
oppure utilizzare la macchina di Turing vista nel teorema dello speedup multinastro (\ref{teorema:speedup_multinastro}) 
per simulare tutti i nastri di $M$ in uno solo, garantendo che $U$ funzioni con 2 nastri. \\

$U$ codifica un nastro di $M$ nella forma $[w]q[u]$ (lo stato fa da separatore tra $w$, i simboli fino alla testina, e $u$, i simboli subito dopo). \\
Se $M$ é nello stato $q$, allora $U$ cerca nel primo nastro, dove c'é $\underline{M}$, le transizioni che iniziano per $q$. \\
Trovatane una, legge i simboli sotto le testine dei $k$ nastri per verificare che siano uguali a quelli della transizione 
(se usiamo $k$ nastri, basta spostare tutte le testine a sinistra di una posizione, altrimenti deve scandire il nastro che contiene la configurazione di $M$). \\
Se sono uguali, allora $U$ applica la transizione. \\
Altrimenti, cerca la prossima quintupla. \\
In caso di errori, la macchina si sposta a destra all'infinito. \\

% \subsubsection*{}
Va da sé che $|x| < N$, e quindi, per qualsiasi funzione propria, $f(|x|) = O(f(N))$.

\subsection{Linguaggio $H_f$}
\[H_f = \{\underline{M}, x \mid M \text{ accetta } x \text{ entro } f(|x|) \text{ passi} \} \]
Questo linguaggio é simile all'halting problem, ma con un limite di tempo che garantisce la terminazione della macchina che lo calcola, e quindi la decidibilitá 
(in pratica, l'halting problem é l'unione di tutti i linguaggi $H_f$.)

\subsection{Lemma 1: Limite superiore alla complessitá di $H_f$}
\[H_f \in \TIME((f(N))^3)\]
\begin{proof}
Utilizziamo la macchina di Turing universale $U$ con input $X = (\underline{M}, x)$, usando un solo nastro per le configurazioni di $M$,
ma aggiungendone uno per separare $x$ da $\underline{M}$ 
e un altro per avere un contatore che garantisca la terminazione 
in $f(|x|)$ transizioni simulate di $M$. \\

$f$ é propria: iniziamo ad applicare $M_f$ (sezione \ref{sec:funzioni_proprie}) per scrivere $\sqcap^{f(|x|)}$ su un nastro, 
in modo da utilizzarlo come contatore del numero di passi durante la simulazione di $M$. \\
Tempo: $O(|x| + f(|x|)) = O(f(|x|)) = O(f(N))$. \\

Poi copia $\underline{M}$ su un altro nastro, cancellandola dal primo e lasciando solo $\triangleright x$. \\
Controlla anche che $\underline{M}$ sia ben formata. \\
Tempo: $O(N)$ per la copia di $\underline{M}$ e la traslazione di $x$. \\

Per simulare $M$, utilizziamo la macchina di Turing vista nel teorema dello speedup multinastro (\ref{teorema:speedup_multinastro}) 
per simulare tutti i nastri di $M$ in uno solo, garantendo che $U$ funzioni con 4 nastri. \\

Se $M$ si trova nello stato $q$, $U$ comincia cercando sul nastro dove ha copiato $\underline{M}$ la codifica di una transizione che inizi con $\underline{q}$.\\
Tempo per cercare la transizione: $O(|\Delta|) = O((\log(N))^2)$. \\
Trovatane una, legge la configurazione di $M$, verificando che i simboli sotto le testine corrispondano a quelli richiesti dalla transizione. \\
Se non corrispondono, deve cercare la prossima transizione e ripetere. \\
La lunghezza del nastro su cui legge la configurazione di $M$ é $k \cdot \ell \cdot O(f(|x|))$.
Di conseguenza, il tempo per leggere la configurazione é $O(f(|x|)) = O(f(N))$.\\
Tempo complessivo per trovare la transizione da applicare: $O(f(N) \cdot (\log(N))^2) = O((f(N))^2)$.\\

Dopo ogni transizione, sposta a destra la testina sul contatore. 
Se il contatore termina, $U$ rigetta.
In questo modo la terminazione é garantita entro $f(|x|) = O(f(N))$ transizioni.\\

Chiaramente il fattore dominante é il tempo della simulazione: 
per $O(f(N))$ volte, bisogna applicare una transizione in tempo $O((f(N))^2)$. 
Il tempo complessivo é:
\[O((f(N))^3)\]
\end{proof}

